% !TEX root = CoeneDylanBP.tex
%%=============================================================================
%% Implementatie van de Meetpipeline
%%=============================================================================

\chapter{Implementatie van de Meetpipeline}
\label{ch:meetpipeline}

Voortbouwend op het systeemontwerp uit Hoofdstuk~\ref{ch:systeemontwerp} beschrijft dit hoofdstuk hoe de effectieve metingen worden uitgevoerd. Achtereenvolgens wordt ingegaan op het detecteren van het exacte landingsmoment, het bepalen van de bijbehorende positie op het trampolinedoek en de berekening van de vluchttijd op basis van de gedetecteerde contactmomenten.

%----------------------------------------------------------------------

\section{Landingsmomentdetectie}
\label{sec:landingsmoment}

Om de landingspositie te kunnen bepalen, moet eerst het exacte frame worden ge{\"i}dentificeerd waarop de springer in aanraking komt het doek. Deze stap is cruciaal, gezien zowel een te vroeg als een te laat gedetecteerd landingsmoment leidt tot een foutieve positieschatting.

\subsection{Initi{\"e}le aanpak: Snelheidsgebaseerde Detectie via Pose Estimation}
\label{subsec:landingsmoment-pose}

De eerste strategie om het landingsmoment te bepalen was gebaseerd op de verticale snelheid van de springer doorheen de sprong. Via een Pose Estimation model werd het zwaartepunt van de springer per frame bepaald. Wanneer de verticale component van dat zwaartepunt abrupt afnam, een teken dat de opwaartse beweging stopt, zou dat het moment van doekcontact aangeven.

In de praktijk bleek deze aanpak echter onbetrouwbaar. YOLO Pose Estimation modellen genereren keypoint-predicties op basis van individuele frames, zonder gebruik te maken van temporele informatie. Bij de snelle bewegingen die kenmerkend zijn voor trampolinespringen veroorzaakt de hoge snelheid van de springer op het moment vlak voor landing een sterke bewegingsonscherpte (\textit{motion blur}) in het beeld. Precies in deze kritieke frames vielen de keypoint-predities frequent weg of werden ze onnauwkeurig, waardoor het snelheidsignaal te onbetrouwbaar was om als detectiemechanisme te dienen.

\subsection{Gekozen aanpak: Binaire Classifier}
\label{subsec:landingsmoment-classifier}

Omdat het snelheidsgebaseerde systeem te sterk afhankelijk was van betrouwbare keypoints, werd overgestapt op een aanpak die het landingsmoment direct op beeldniveau detecteert. Hiervoor werd een YOLO11n classificatiemodel \autocite{Jocher_Ultralytics_YOLO_2023} getraind als binaire classifier met twee klassen: \textit{Landed} en \textit{Not Landed}.

Als invoer ontvangt de classifier per frame een uitsnede van het trampolinebed, bepaald als bijproduct van de hoekpuntendetectie uit Sectie~\ref{sec:trampolinedetectie}. Door de invoer te beperken tot de trampoline-regio wordt de achtergrond en de omgeving van het beeld buiten beschouwing gelaten, wat de focus van het model op de relevante informatie vergroot. Het model leert zo om direct uit het visuele patroon van het doek, de aanwezigheid van een springer, de vervorming van het doek en de schaduwpatronen te bepalen of er contact is.

Het eerste opeenvolgende frame waarop de classifier de klasse \textit{Landed} toekent, wordt gebruikt als het landingsframe. Dit frame dient vervolgens als invoer voor de landingspositiebepaling (Sectie~\ref{sec:landingspositie}) en als ankerpunt voor de ToF-berekening (Sectie~\ref{sec:tof}).

Ondanks de relatief beperkte omvang van de trainingsset presteerde het model opvallend goed. De eenvoud van de classificatietaak, het onderscheid tussen een leeg doek en een doek met een springer, maakt het probleem voldoende afgebakend om met weinig data goed te generaliseren.

%----------------------------------------------------------------------

\section{Landingspositiebepaling}
\label{sec:landingspositie}

De landingspositie van de springer wordt bepaald op het gedetecteerde landingsframe (zie Sectie~\ref{sec:landingsmoment}). Het bepalen van een betrouwbaar referentiepunt voor deze positie bleek iteratief: meerdere aanpakken werden achtereenvolgens onderzocht en verfijnd.

\subsection{Poging 1: YOLO-Pose Enkelmiddelpunt}
\label{subsec:landingspositie-yolo}

De eerste aanpak maakte gebruik van een YOLO11l-pose model \autocite{Jocher_Ultralytics_YOLO_2023} om de lichaamspositie van de springer te bepalen. In het COCO-keypointformaat zijn de enkels gedefinieerd als keypoints 15 (linkerenkel) en 16 (rechterenkel). Het rekenkundig middelpunt tussen deze twee keypoints werd gebruikt als referentiepunt voor de landingspositie, omdat dit beter aansluit bij het werkelijke contactpunt van de voeten dan het zwaartepunt van de volledige bounding box.

In de praktijk deden zich echter twee fundamentele problemen voor. Ten eerste veroorzaakte de hoge snelheid van de springer op het landingsmoment sterke bewegingswaas in het beeld. Omdat YOLO-pose elk frame onafhankelijk verwerkt zonder temporele context, resulteerde dit in frequente uitval van de enkelpredities op precies de frames die het meest relevant waren. Ten tweede bevonden zich in het wedstrijdmateriaal regelmatig \textit{spotters} --- begeleiders die vlakbij de trampoline staan om de springer bij te staan. Deze personen werden soms verkeerdelijk als de springer ge{\"i}dentificeerd, wat leidde tot volledig foutieve posities.

\subsection{Poging 2: Lineaire Interpolatie}
\label{subsec:landingspositie-interpolatie}

Om de hiaten door missende keypointpredities op te vangen, werd in een tweede iteratie lineaire interpolatie toegepast. Tussen twee frames met een geldige detectie werd de positie van de ontbrekende frames geschat als een lineaire overgang tussen de twee gekende posities.

Deze aanpak verbeterde de continuïteit van de positiereeks maar leverde onvoldoende nauwkeurige resultaten op. De beweging van een springend gymnast is fundamenteel niet-lineair: na afstoot versnelt de springer, vertraagt bij het hoogste punt, en versnelt opnieuw richting het doek. Een lineaire interpolatie gaat voorbij aan deze dynamiek. Bovendien waren de aaneengesloten periodes zonder geldige detectie vaak lang, waardoor de ge{\"i}nterpoleerde waarden ver afweken van de werkelijke positie.

\subsection{Poging 3: ViTPose}
\label{subsec:landingspositie-vitpose}

Als robuuster alternatief voor YOLO-pose werd ViTPose \autocite{Xu2022} ge{\"e}valueerd. Dit model, toegelicht in Sectie~\ref{subsec:pose-estimation} van de literatuurstudie, maakt gebruik van een Vision Transformer-backbone die door zijn globale attentiemechanisme beter in staat is om complexe ruimtelijke patronen te herkennen, ook bij deels versluierde of vervormde lichamen.

ViTPose leverde inderdaad betere keypointpredities op dan YOLO-pose, met name in frames met matige bewegingswaas. Bij de snelste bewegingen --- de frames net voor en op het moment van landing --- bleven de predities echter onvoldoende betrouwbaar. De fundamentele beperking van het werken met enkelvoudige frames zonder temporele context gold voor ViTPose evenzeer als voor YOLO-pose. Dit leidde tot de beslissing om de volledige keypoint-gebaseerde aanpak te verlaten en over te stappen op segmentatie.

\subsection{Finale aanpak: SAM-2 Segmentatie}
\label{subsec:landingspositie-sam2}

In plaats van te proberen de positie van de springer te reconstrueren via expliciete lichaamspunten, werd de probleemstelling verlegd: in plaats van \textit{waar zijn de enkels?} werd de vraag \textit{welk pixel in het beeld staat het laagst en behoort tot de springer?}. Dit leidde tot een segmentatiebenadering op basis van SAM-2 \autocite{Ravi2024}, toegelicht in Sectie~\ref{subsec:sam2}.

De pipeline voor de finale aanpak verloopt als volgt:
\begin{enumerate}
    \item Een standaard YOLO11n detectormodel lokaliseert de springer in het landingsframe en geeft een bounding box terug.
    \item Deze bounding box wordt als \textit{prompt} meegegeven aan SAM-2-tiny, dat op basis hiervan een pixel-nauwkeurig segmentatiemasker van de springer genereert.
    \item De laagste pixel van dit masker (het pixel met de grootste $y$-co{\"o}rdinaat in beeldco{\"o}rdinaten) wordt ge{\"i}dentificeerd als het voetcontactpunt.
    \item Dit contactpunt wordt via de homografiematrix $H$ (Sectie~\ref{sec:homografie}) getransformeerd naar een metrische positie op het trampolinebed.
\end{enumerate}

Het cruciale voordeel van deze aanpak is dat er geen expliciete keypointpredities meer nodig zijn. SAM-2 genereert zijn segmentatiemasker op basis van de globale beeldstructuur, waardoor het model robuuster is tegen de bewegingswaas die de pose-gebaseerde methoden parten speelde. Bovendien is de aanpak minder gevoelig voor het probleem van spotters: de bounding box van YOLO bepaalt van welke persoon het masker wordt gemaakt, en door bewust de persoon op het trampolinedoek te selecteren kan de verwarring met begeleiders worden vermeden.

Dit berekende contactpunt in pixelco{\"o}rdinaten wordt vervolgens met behulp van de homografiematrix $H$ (Sectie~\ref{sec:homografie}) getransformeerd naar het co{\"o}rdinaat in het stelsel van de trampoline. Op basis van deze transformatie wordt bepaald in welke HD-zone de landing valt.

%----------------------------------------------------------------------

\section{ToF-schatting}
\label{sec:tof}

Naast de horizontale verplaatsing is de vluchttijd (ToF) een tweede kerncomponent van de eindscore. Het doel van deze sectie is om de ToF te bepalen op basis van de gedetecteerde landingsmomenten en de framerate van de inputvideo, zonder dat er extra sensoren aan te pas komen.

Omdat de video's vanuit een zijdelings perspectief zijn gefilmd, is een directe schatting van de absolute vluchthoogte ($z$-co{\"o}rdinaat) niet haalbaar zonder dieptesensor. In plaats daarvan wordt de ToF bepaald via \textit{contactdetectie}: de landingsframes gedetecteerd in Sectie~\ref{sec:landingsmoment} fungeren als ankerpunten voor de tijdsmeting.

De ToF-berekening verloopt als volgt: voor elke sprong in een reeks wordt het frame-verschil bepaald tussen het landingsframe van de vorige sprong en het landingsframe van de volgende sprong. Dit verschil, gedeeld door de framerate van de video (frames per seconde), geeft de duur van één vluchttijdsinterval. De totale ToF-score is de som van alle vluchttijdsintervallen over de volledige reeks van tien sprongen.

De nauwkeurigheid van deze methode is inherent begrensd door de framerate van de video. Elke landing kan slechts met een precisie van $\tfrac{1}{\text{fps}}$ seconden worden bepaald. Bij een framerate van 25 fps bedraagt deze onzekerheid 40\,ms per landing, wat over een volledige reeks kan oplopen tot een cumulatieve fout van honderden milliseconden. Dit wordt verder besproken in Hoofdstuk~\ref{ch:evaluatie}.
